\begin{frame}{¿Cómo verificar la equidistribución?}

\begin{block}{Problema}
La definición exige comprobar que
\[
\lim_{N\to\infty}\frac{A_N((a,b),x)}{N}=b-a
\]
para todo intervalo $(a,b)\subset\mathbb T$.
\end{block}

\pause

\begin{alertblock}{Objetivo}
Buscar criterios más prácticos para demostrar
que una sucesión es equidistribuida.
\end{alertblock}

\vspace{0.3cm}

En este capítulo presentamos dos herramientas:
\begin{enumerate}
\item Caracterización funcional.
\item Criterio de Weyl.
\end{enumerate}

\end{frame}

% Caracterización funcional
\begin{frame}{Caracterización funcional}

\begin{block}{Teorema}
Las siguientes afirmaciones son equivalentes:
\begin{enumerate}
\item $(x_n)$ es equidistribuida módulo $1$.
\item Para toda función Riemann integrable
$f:\mathbb T\to\mathbb R$,
\[
\lim_{N\to\infty}\frac{1}{N}\sum_{n=1}^{N}f(x_n)=\int_0^1 f(x)\,dx.
\]
\end{enumerate}
\end{block}

\begin{alertblock}{Interpretación}
La media discreta sobre la sucesión coincide
con la media continua sobre el toro.
\end{alertblock}

\end{frame}

% Idea de geometrica
\begin{frame}{Interpretación geométrica}

\[
\frac1N\sum_{n=1}^{N}f(x_n)
\]

aproxima

\[
\int_0^1 f(x)\,dx.
\]

\vspace{0.4cm}

\begin{center}
% \includegraphics[width=.7\textwidth]{media_integral.pdf}
\end{center}

\vspace{0.2cm}

\begin{alertblock}{Idea}
Una sucesión equidistribuida se comporta como una
muestra uniforme del intervalo $[0,1)$.
\end{alertblock}

\end{frame}

% Criterio de Weyl
\begin{frame}{Criterio de Weyl}

\begin{block}{Teorema}
Las siguientes afirmaciones son equivalentes:

\begin{enumerate}
\item $(x_n)$ es equidistribuida módulo $1$.
\item Para todo $k\in\mathbb Z\setminus\{0\}$,

\[
\frac1N
\sum_{n=1}^{N}
e^{2\pi i k x_n}
\longrightarrow 0.
\]

\end{enumerate}
\end{block}

\pause

\begin{alertblock}{Ventaja}
El problema se reduce al estudio de sumas
de exponenciales complejas.
\end{alertblock}

\pause

\[
\text{Equidistribución}
\quad \Longrightarrow \quad
\text{Análisis de Fourier}
\]

\end{frame}

% Aplicación del criterio de Weyl
\begin{frame}{Aplicación: la sucesión $n\alpha$}

\begin{block}{Resultado}
Si $\alpha\notin\mathbb Q$, entonces $(n\alpha)_{n\geq1}$
es equidistribuida módulo $1$.
\end{block}

Por el criterio de Weyl basta estudiar
\[\frac{1}{N}\sum_{n=1}^{N}e^{2\pi i k n\alpha}.\]

Esta suma es geométrica y satisface
\[
\left|\frac{1}{N}\sum_{n=1}^{N} e^{2\pi i k n\alpha}\right|
\le\frac{2}{N|e^{2\pi i k\alpha}-1|}.
\]

Por tanto,
\[\lim_{N\to\infty}\frac{1}{N}\sum_{n=1}^{N}e^{2\pi i k n\alpha}=0.\]
\end{frame}